\section{模型改进与推广}

\subsection{接入真实流程时间}

若能获得预约时刻、到达时刻、扫查开始与结束、报告提交和审核时刻，可将当前单一“最后任务结束”拆为
\[
\text{开单}\rightarrow\text{预约}\rightarrow\text{到达}\rightarrow\text{扫查}\rightarrow\text{报告}\rightarrow\text{审核}.
\]
用生存分析或分位数回归估计等待、服务和报告延迟分布，再将48小时完成写成机会约束
\[
\Pr\{C_i^{\mathrm{report}}\le p_i+48\unit{h}\}\ge1-\epsilon_i.
\]
这样可以区分预约拥堵和报告流程延迟，也能对不同项目给出具有统计覆盖率的缓冲，而不是统一取上界。

\subsection{医生实名资质与排班联动}

获得未来医生集合$\mathcal A$、医生—项目资质矩阵$h_{aj}$和实名班次$B_{at}$后，引入医生指派变量$z_{ijat}$：
\[
z_{ijat}\le h_{aj},\qquad
\sum_{i,j}z_{ijat}\le B_{at},\qquad
\sum_a z_{ijat}=\sum_r x_{ijrt}.
\]
这会把当前总并发近似升级为设备—医生—项目三维匹配。若同一医生在一段连续时间内尽量固定设备，还可加入换房惩罚，减少人员走动和消毒切换。

\subsection{滚动时域与在线重排}

全年回放用于评价，而真实运行应采用滚动时域。每天或每小时重建集合：
\[
\I^{(k)}=\I_{\mathrm{unfinished}}^{(k)}\cup
\I_{\mathrm{booked}}^{(k:k+H)}\cup
\I_{\mathrm{new}}^{(k)}.
\]
固定已经开始或即将开始的项目，只对窗口$H$内其余事件重排。目标中加入改约惩罚
\[
\lambda_{\mathrm{chg}}\sum_i|s_i^{(k)}-s_i^{(k-1)}|,
\]
在提高完成率的同时控制患者通知成本。对于床旁急诊或设备故障，可局部释放受影响任务并重算，不必推翻全年方案。

\subsection{公平与服务等级扩展}

对科室$g$、多项目组和高严重度组定义分组完成率$R_g$，增加
\[
R_g\ge \rho_g,\qquad
\max_g R_g-\min_gR_g\le\Delta.
\]
也可采用分层词典序目标：先满足急危重和政策下限，再最大化住院总完成，随后优化门诊、体检和等待。相较把所有患者压成单一权重，这种形式更便于医院解释。

\subsection{容量建设情景}

校准链表明项目专项能力和三类背景竞争是主要瓶颈。可把设备增购、工作时段延长和医生增班设为上层决策：
\[
\min\ \sum_r c_r^{\mathrm{device}}u_r+\sum_t c_t^{\mathrm{overtime}}v_t
+\kappa(1-R_I),
\]
其中$u_r$表示新增某类设备，$v_t$表示扩展时段。下层仍为本文患者级排程。通过双层或样本平均近似，可比较“增加普通设备”“增加床旁设备”“增加产科/心脏专项设备”和“增加医生班次”的边际收益。由于C1--C3已显示总医生并发不是纯住院首要瓶颈，优先增购普通同质设备未必有效，必须与专项能力共同评估。

\subsection{向其他医疗预约问题推广}

该框架可推广到MRI、CT、内镜、放疗和手术室预约：患者事件对应诊疗申请，项目任务对应序列，设备兼容对应检查协议或手术间资质，医生容量对应技师、麻醉和术者，准备条件对应禁食、造影和消毒。推广时应保留三项原则：以患者完整流程为目标；能力集合必须由真实资质产生；统计参数和评价期严格隔离。迁移到新场景时，资源数量、资质集合、工作时段、准备和转运时间必须重新估计，不能直接复制本题参数。
